\documentclass[11pt]{article}

\usepackage{multicol}
\usepackage{times}
\usepackage{graphicx}
\usepackage{epsfig}

\textheight 9in
\textwidth 6.5in
\oddsidemargin 0in
\evensidemargin 0in
\topmargin -.5in
\pagestyle{empty}
\newcommand{\setspace}[1]{\renewcommand{\baselinestretch}{#1}\large\normalsize}
\setspace{1.1}
\renewcommand{\textfraction}{0.1}
\renewcommand{\topfraction}{0.9}
\newcommand{\tods}{{\em TODS\/}}

\begin{document}
\begin{center}
{\Large\bf CMM and {\em TODS}}\\[2ex]
{\large\it  Richard Snodgrass}\\[1ex]
{\tt rts@cs.arizona.edu}\\[3ex]
\end{center}

The {\em Capability Maturity Model}~\cite{CMM} is an orderly way for
organizations to determine the capabilities of their current processes for
developing software and to establish priorities for improvement~\cite{Humphrey}. It
defines five levels of progressively more mature process
capability~\cite{Paulk}.

\begin{description}
\item[Level 1: Initial] The software process is characterized as ad hoc, and
  occasionally even chaotic. Few processes are defined, and success depends
  on individual effort.

\item[Level 2: Repeatable] Basic project management processes are
  established to track cost, schedule, and functionality. The necessary
  process discipline is in place to repeat earlier successes on projects
  with similar applications.

\item[Level 3: Defined] The software process for both management and
  engineering activities is documented, standardized, and integrated into an
  organization-wide software process. All projects use a documented and
  approved version of the organization's process for developing and
  maintaining software. This level includes all the characteristics defined
  for level 2.

\item[Level 4: Managed] Detailed measures of the software process and
  product quality are collected. Both the software process and products are
  quantitatively understood and controlled using detailed measures. This
  level includes all the characteristics defined for level 3.

\item[Level 5: Optimizing] Continuous process improvement is enabled by
  quantitative feedback from the process and from testing innovative ideas
  and technologies. This level includes all the characteristics defined
  for level 4.
\end{description}

You may be asking, what does a maturity model for software development have
to do with databases generally and with \tods\ in particular? Well, CMM has
been applied to personnel management, quality management, and even weapons
system development. And it can be used as a framework for evaluating the
journal review process, as we will do here.

Manuscript review at \tods\ started, logically, at Level 1. In 2001, the ACM
Publications Board approved a broad
policy~\cite{ACMPubs,SnodgrassCACM,SnodgrassSIGMOD} that raised publishing
of ACM journals and transactions to Level 2. In 2003 ACM adopted the
Manuscript Central\footnote{\tt http://acm.manuscriptcentral.com} web-based
manuscript tracking system~\cite{SnodgrassMC}, raising its manuscript
reviewing process to Level 3.

In parallel with these efforts at the ACM Publications Board level, I have
been refining the reviewing process for \tods. In October 2003 I released
the first edition of the {\bf {\em ACM TODS} Associate Editor Manual}, with
revisions in April 2004 and October 2004. This manual, at 22 pages, is quite
detailed. 

I have also been collecting detailed statistics since
July 2001. Some of these statistics are
reported on the \tods\ web site\footnote{\tt
http://www.acm.org/tods/TurnaroundTime.html}: turnaround time, article
length, number of articles, and end-to-end time~\cite{SnodgrassSIGMOD}. I
have also kept records on the turnaround time of individual Associate
Editors, and have closely monitored the progress of individual papers.

Through these efforts, and through a series of internal policies regarding
the reviewing process that has been adopted by the \tods\ Editorial Board,
all of the statistics has improved, some
considerably~\cite{SnodgrassSept03}. Average turnaround time is now down to
13 weeks, average article length has been brought down to levels last seen
in the mid-1990s (under 40 pages), the number of articles per volume is back
up to that last experienced in the early 1990's (21 articles per year), and
average end-to-end time is down to 17 months, last experienced in the
1970's.

The result is that \tods\ is now operating at CMM Level 5.

Why should you, dear reader, care about internal processes at \tods?
The short answer is that by being at Level 5, \tods\ can provide assurances
as to how {\em your} submission will be handled.

Average turnaround and end-to-end times are nice, but what authors really
care about is how long {\em their} submission will take to be
reviewed. Addressing this concern involves both average and {\em maximum}
times. A low average turnaround time is of little reassurance to someone
experiencing an abnormally long turnaround time. As an example, while the
average turnaround time for papers submitted in January 2002 to \tods\ was a
quite reasonable 5.5 months, one paper submitted that month had to wait
almost nine (!) months for a decision.

By virtue of being at CMM Level 5, the variance of the turnaround time could
be monitored and improved, as shown in Figure~\ref{fig:turnaround}.

The turnaround time has been slowly decreasing over the past four years.
This figure shows four sets of data. The bottom line is the {\em average
turnaround time}, a moving average of the turnaround time for papers
submitted in the indicated month. To smooth monthly variations, the moving
average includes all of the submissions for the previous year. Each data
point represents dozens of papers. The value for January 2005, 12.5 weeks,
is the average turnaround time for all of the papers submitted between
(inclusive) February 2004 and January 2005.

The next line up is the average turnaround time for external reviews only, a
moving average of the turnaround time for papers submitted in the indicated
month. This includes only submissions that went out to external reviewers
and specifically excludes desk rejects. The value for January
2005, \hbox{15.6~weeks}, is the average turnaround time for external reviews
of all the papers submitted during the year up through January 2005.

The points, one per month, denote the maximum or peak turnaround time for
submissions in the indicated month. Each point represents a single,
unusually slow paper submitted during the indicated month. For all the
papers submitted in January 2005, the longest turnaround time was 4.9
months (21 weeks).

In terms of turnaround time, \tods\ at 12.5 weeks is now equivalent to
conferences (as exemplified by \hbox{SIGMOD} and PODS at 12 weeks), while
being more flexible in not imposing a submission deadline.

The straight line is the {\em committed maximum turnaround time}, the
boundary that the Editorial Board has committed to not exceed, for any
submission. Several years ago the Editorial Board established a formal
policy stating its commitment to providing an editorial decision within 6
months~\cite{SnodgrassSept03}. \tods\ thus joined conferences in
guaranteeing a stated turnaround time.

Due to the rigorous application of CMM Level 5, of continuous process
improvement as exemplified by the steady lowering of average turnaround time
and the compression of the variance in turnaround time by a factor of two,
I can announce that the Editorial Board is now committed to providing an
editorial decision within {\em five} months, starting with submissions in
2004. As depicted in the figure., we have met this stated commitment for the
past thirteen months. As of the writing of this column (June 29, 2005),
all manuscripts submitted before February 1 of this year have been processed
and editorial decisions rendered.

\begin{figure}[t]
\begin{center}
\epsfig{figure=sep05tt.eps,width=4.5in}
\end{center}
\caption{ACM {\em TODS} Turnaround Time\label{fig:turnaround}}
\end{figure}

That \tods\ now matches conferences in terms of turnaround time is a
testament to the hard work of two groups of people: reviewers and the
editorial board. I will recognize the reviewers in a future column, but here
I wish to thank the following people, who comprise the \tods\ Editorial
Board, for their dedicated effort work in achieving very fast decisions
while upholding very high standards.

\vspace{4ex}

\begin{multicols}{2}
\noindent
Surajit Chaudhuri, Microsoft Research\\
Jan Chomicki, SUNY Buffalo\\
Mary Fernandez, AT\&T Labs\\
Michael Franklin, Univ.~of California at Berkeley\\
Luis Gravano, Columbia University\\
Ralf Hartmut G\"{u}ting, Fernuniversit\"{a}t Hagen\\
Richard Hull, Bell Labs\\
Christian S. Jensen, Aalborg University\\
Hank Korth, Lehigh University\\
\\
Donald Kossmann, ETH Zurich\\
Heikki Mannila, University of Helsinki\\
Z. Meral \"Ozsoyo\v{g}lu, Case Western Reserve\\
Raghu Ramakrishnan, University of Wisconsin
Arnie Rosenthal, MITRE\\
Betty Salzberg, Northeastern University\\
Sunita Sarawagi, IIT Bombay\\
Dan Suciu, University of Washington\\
Jennifer Widom, Stanford University\\
\end{multicols}

\noindent
These 18 people are providing a truly valuable service to readers, to
authors, and to reviewers. When you see these people, please thank them
personally for their role in achieving quick reviews of submitted papers.

\pagebreak

\begin{thebibliography}{[10]}

\bibitem{ACMPubs}
ACM Publications Board, ``Rights and Responsibilities in ACM Publishing,''
approved June 27, 2001. ({\tt http://www.acm.org/pubs/rights.html})

\bibitem{Humphrey}
Watts S.~Humphrey, {\bf A Discipline for Software Engineering},
Addison-Wesley, 1995.

\bibitem{Paulk}
M.~C.~Paulk, Bill Curtis, and M.~B.~Chrisis, ``Capability Maturity Model for
Software, Version 1.1,'' Software Engineering Institute Technical Report,
CMU/SEI-93-TR, February 24, 1993.

\bibitem{CMM}
SEI, {\bf The Capability Maturity Model: Guidelines for Improving the Software
Process}, Software Engineering Inst. Carnegie Mellon Univ., Addison-Wesley,
1995.

\bibitem{SnodgrassCACM}
Richard T.~Snodgrass, ``Rights and Responsibilities in ACM Publishing,''
{\em CACM} 45(2): 97--101, February 2002.

\bibitem{SnodgrassSIGMOD}
Richard T.~Snodgrass, ``Rights of {\em TODS} Readers, Authors and
Reviewers,'' {\em SIGMOD Record}, 31(4):5--9, December 2002.

\bibitem{SnodgrassMC}
Richard T.~Snodgrass, ``ACM {\em TODS} in this Internet Age,'' {\em
  SIGMOD Record}, 32(1):4--5, March 2003.

\bibitem{SnodgrassSept03}
Richard T.~Snodgrass, ``Journal Relevance,'' {\em SIGMOD Record},
31(3):11--15, September 2003.
\end{thebibliography}
\end{document}
